% TOP_PROBLEM: {"id": 2800802, "title": "10 Lectures and 42 Open Problems — Sum of Squares approximation ratio for Max-Cut", "category_id": 15, "category": {"id": 15, "name": "computer_science", "display_name": "Computer Science", "description": "Computational complexity, algorithms, and theoretical CS.", "slug": "computer-science", "order_index": 15, "created_at": "2026-07-31T15:26:25.671Z"}, "status": "open", "problem_number": "AMR-027-0802", "set_id": 13}
% FIRST_POSTED: 2026-10-02
% ENTRY_KIND: statement_only
% UnsolvedMath ID 2800802; source catalogue code AMR-027-0802
% !TeX program = lualatex
% Definition and sources only; this file is not an LLM solution attempt.
\documentclass[11pt,a4paper]{article}
\usepackage[margin=25mm]{geometry}
\usepackage{fontspec}
\setmainfont{lmroman10-regular.otf}[BoldFont=lmroman10-bold.otf,
 ItalicFont=lmroman10-italic.otf,BoldItalicFont=lmroman10-bolditalic.otf]
\usepackage{amsmath,amssymb,mathtools}
\usepackage{unicode-math}
\setmathfont{latinmodern-math.otf}
\AtBeginDocument{\let\setminus\smallsetminus}
\usepackage{xurl,hyperref}
\hypersetup{colorlinks=true,urlcolor=blue}
\setcounter{secnumdepth}{0}
\setlength{\parindent}{0pt}
\setlength{\parskip}{5pt}
\setlength{\emergencystretch}{3em}
\begin{document}
\section*{10 Lectures and 42 Open Problems — Sum of Squares approximation ratio for Max-Cut}\label{2800802}
{\small UnsolvedMath ID 2800802 · AMR-027-0802 · Computer Science · AMR Open Problem Lists}
\subsection{Definitions and mathematical statement}
Let $G=(V,E,w)$ be a finite graph with nonnegative rational edge weights
of positive total sum. For signs $x_v\in\{-1,1\}$, its cut weight is
$f_G(x)=\sum_{uv\in E}w_{uv}(1-x_ux_v)/2$; write
$\operatorname{OPT}(G)=\max_x f_G(x)$.

A degree-four pseudoexpectation is a linear functional $L$ on real
polynomials of total degree at most four satisfying
\[
 L(1)=1,\qquad L(q^2)\ge0\quad(\deg q\le2),\qquad
 L((x_v^2-1)q)=0\quad(\deg q\le2, v\in V).
\]
Define $\operatorname{SOS}_4(G)=\max_L L(f_G)$. These conditions are
a semidefinite relaxation of expectations over actual cuts, so
$\operatorname{SOS}_4(G)\ge\operatorname{OPT}(G)$.

Determine its worst-case integrality ratio
\[
       \alpha_4=\inf_G
             \frac{\operatorname{OPT}(G)}{\operatorname{SOS}_4(G)},
\]
and the best polynomial-time rounding guarantee obtainable from this
relaxation. The ratio convention is stated explicitly: larger values
mean a tighter relaxation. An algorithm attaining a proposed bound
and examples obstructing improvement play different roles.

The source also asks about any other fixed even degree $2r$, obtained
by replacing the degree bounds by $2r$, $r$, and $2r-2$, respectively.
The degree is fixed independently of $|V|$. In the near-bipartite
regime, where a cut contains a $1-\varepsilon$ fraction of total
weight, can degree four improve the degree-two guarantee
$1-O(\sqrt\varepsilon)$? Merely adding constraints and changing some
instances' optimum does not determine the worst-case ratio.
\subsection{Short English statement}
Find the worst-case strength and rounding guarantees of degree-four, and other fixed-degree, sum-of-squares relaxations for Max-Cut.
\subsection{Sources}
\begin{itemize}
\item[S1] ULAM.AI and the UnsolvedMath Contributors, supplied problems.json, record 2800802 (AMR-027-0802), AMR Open Problem Lists. Catalogue identity and source transcription; CC BY 4.0.
\url{https://huggingface.co/datasets/ulamai/UnsolvedMath}.
\item[S2] Bandeira - 42 Open Problems in Data Science (2016), Open Problem 8.2. Original source indexed by AMR-027-0802.
\url{https://afonsobandeira.wordpress.com/2015/12/04/10l42pmaxcut/}.
\item[S3] A. S. Bandeira, Ten Lectures and Forty-Two Open Problems in the Mathematics of Data Science, corresponding numbered problem and surrounding definitions.
\url{https://people.math.ethz.ch/~abandeira/TenLecturesFortyTwoProblems.pdf}.
\end{itemize}
\end{document}
